% !TeX root = ../nc-quadric-avatar.tex
% 2026-11-23 / v4.25。可逆cubicとmass壁の相対QP群oid。

\subsection{八矢QPから有限cubic foundationへの全相対群oid}
\label{subsec:qa-cubic-qp-foundation-groupoid}

本項では前記の四頂点とweightを保持し，標数零の体$k$上で扱う．
以下の有限Grassmann商の同型自体は任意の基礎体上で成立する．
八矢箙$Q_8$の隣接矢束を
\[
 U:0\to1,\qquad V:1\to2,\qquad W:2\to3,\qquad T:3\to0
\]
（いずれも階数二）とし，
\(\mathcal V=W\otimes V\otimes U\) と置く．
頂点で始点を固定した巡回道を用いれば，
weight四のpotentialは正確に
\[
 H\in\mathcal V\otimes T,\qquad
 \partial_T H:\ T^\vee\longrightarrow\mathcal V
 \tag{QP8}\label{eq:qa-cubic-qp-flattening}
\]
に対応する．ここでは巡回同値を取っても余分な係数商は生じない．
任意のスキーム上で \(\partial_T H\) が階数二の局所自由直和因子
である開を \(P_8^{(2)}\) と書く．
これはrank-two flatteningの全minorの単位イデアル条件である．
\[
 G_8=\GL(U)\times\GL(V)\times\GL(W)\times\GL(T),
 \qquad G_3=\GL(U)\times\GL(V)\times\GL(W)
\]
を重み付きright-equivalence群とする．$Q_8$では任意の
重み一の矢はその始点終点の二矢の線型結合であり，
同じ始点終点の重み一の長い道は存在しない．従って
重みを保つ右同値の全群は正に$G_8$である．
非重み付き完備道代数の全自己同型群と同一視しない．

\begin{theorem}[八矢QPと階数二関係束の相対groupoid同値]
\label{thm:qa-cubic-qp-grassmann-groupoid}
次のfppfスタックの自然な同値がある：
\[
 \boxed{\quad
 [P_8^{(2)}/G_8]\ \simeq\
 [\operatorname{Gr}(2,\mathcal V)/G_3]\quad}.
 \tag{QG}\label{eq:qa-cubic-qp-grassmann-stack}
\]
左辺はweight四の四次potentialと重み付きright-equivalenceを
すべての同型として持つ．右辺は三つの階数二束と
$\mathcal V$内の階数二直和因子$R$を同型とともに持つ．
同値は$H$を
$R=\operatorname{im}(\partial_T H)$へ送る．
任意のスキームへの基底変更，非被約基底，Isom-sheafおよび
慣性群スキームと整合する．
\end{theorem}
\begin{proof}
$G_3$をいったん忘れ，各矢束を固定する．
$\operatorname{Gr}(2,\mathcal V)$の普遍階数二部分束を
$\mathcal R\subset\mathcal V$とする．
$\partial_T H$の像が$\mathcal R$であるという条件は，
$T^\vee\xrightarrow{\sim}\mathcal R$という同型を選ぶことに等しい．
従って
\[
 P_8^{(2)}
 =\operatorname{Isom}_{\operatorname{Gr}(2,\mathcal V)}
       (T^\vee,\mathcal R)
\]
は普遍部分束の$\GL(T)$-torsorであり，Zariski局所に自明である．
$\GL(T)$は$P_8^{(2)}$上自由に作用し，
その幾何商およびfppf商は正確にGrassmannスキームとなる．
このtorsorは$G_3$-同変なので，残りの群による商を取れば
表示した二商スタックの同値が得られる．
具体的な逆関手では，$R$に対して$T=R^\vee$と置き，
$\mathcal V\otimes R^\vee$内の恒等テンソルをpotentialとする．
基底変更に伴う第四矢束の自動的な同型は全right-equivalenceを回収する．
単なる幾何点の全単射でなくtorsor同値なので，
相対同型層と慣性も回収する．
\end{proof}

\begin{proposition}[四頂点finite foundationの全積を保つ同型]
\label{prop:qa-cubic-qp-finite-foundation-groupoid}
有限代数スタック
\(\mathfrak F_8^{\rm lin}=[(Z_a\cap Z_b)/G]\)
（命題~\ref{thm:qa-foundation-open-diagram}のunrigidified版）
には自然な同値
\[
 \boxed{\quad
 \mathfrak F_8^{\rm lin}
 \ \simeq\
 [\operatorname{Gr}(2,\mathcal V)/G_3]
 \ \simeq\ [P_8^{(2)}/G_8]\quad}.
 \tag{FG}\label{eq:qa-cubic-qp-finite-groupoid}
\]
ここでは有限代数の階数二十四の全成分・全積を保存する．
QPからの関手は、最後の矢群$T$に沿うcyclic derivativesが張る
三段道の関係束を取り，四頂点のdirected windowの有限End代数を
回収する．Jacobian商の全次数AS正則性は主張しない．
\end{proposition}
\begin{proof}
有限foundationの三つの隣接Hom束を$U,V,W$とする．
$Z_a\cap Z_b$では二つの乗法
\[
 a:V\otimes U\xrightarrow{\sim}E,\qquad
 b:W\otimes V\xrightarrow{\sim}F
\]
が可逆である．
従って$E=V\otimes U$，$F=W\otimes V$へ基底変更で
厳密に規格化できる．
結合律$c(1_W\otimes a)=d(b\otimes1_U)$により
$W\otimes E$と$F\otimes U$の両経路の合成は，同一の写像
\[
 m:\mathcal V=W\otimes V\otimes U\longrightarrow N
\]
となる．元の$(c,d)$の合併全射性は$m$の全射性と同値である．
$\operatorname{rank}\mathcal V=8$，$\operatorname{rank}N=6$だから，
$R=\ker m$は階数二の直和因子である．逆に任意の$R$から
\[
 E=V\otimes U,\quad F=W\otimes V,\quad
 N=\mathcal V/R
\]
および二つの標準的な商写像$c,d$を取れば，全積と結合律を
満たす有限directed代数が得られる．この構成は互いに逆で，
$E,F,N$のbasis torsorも含めて任意の基底変更と可換する．
したがって有限代数の全isomorphism groupoidは
$[\operatorname{Gr}(2,\mathcal V)/G_3]$に等しい．
定理~\ref{thm:qa-cubic-qp-grassmann-groupoid}と合わせれば結論が従う．

さらに八矢QPの微分$\partial_T H$は二つの長さ三の関係を与える．
四頂点のdirected windowでは頂点$3\to0$の矢を通る道は存在しない．
従って他のcyclic derivativesに由来する余分な関係はなく，
最初の四頂点の有限Endの長さ三の商は正確に$\mathcal V/R$である．
二段道には関係がないので，有限代数の全積が一致する．
\end{proof}

\begin{corollary}[頂点スカラーの慣性とcubic $H_2$の相対認識]
\label{cor:qa-cubic-qp-pointed-h2-gerbe}
前命題の$G_3$の三つのscalar部分群
\(K_{\rm vert}\simeq\Gm^3\)はGrassmann上自明に作用する．
その剛体化は有限foundationの
\(\mathfrak F_8=[(Z_a\cap Z_b)/(G/K)]\) と一致し，
\[
 \mathfrak F_8\simeq
 [\operatorname{Gr}(2,\mathcal V)/\operatorname{PGL}_2^3].
 \tag{Fr}\label{eq:qa-cubic-qp-rigid-foundation}
\]
八矢側では，四つの矢群のscalarの積が一となる
\[
 K_{\rm QP}=
 \{(\alpha,\beta,\gamma,\delta)\in\Gm^4:
                 \alpha\beta\gamma\delta=1\}
 \simeq\Gm^3
\]
がpotentialを点ごとに固定する．
$G_8/K_{\rm QP}$への変更が正に上の剛体化を与え，
幾何点の安定化群にはfppf群スキームの完全列
\[
 1\longrightarrow K_{\rm QP}\longrightarrow
 \operatorname{Aut}_{[P_8^{(2)}/G_8]}(H)
 \longrightarrow\operatorname{Aut}_{\mathfrak F_8}(R)
 \longrightarrow1
 \tag{St}\label{eq:qa-cubic-qp-stabilizer-exact}
\]
がある．

四重テンソルのVan den Bergh geometric開$\mathscr G$を
指定すると、この同値はその不変開へ制限される．
幾何開では$\partial_T H$が階数二である．
既存の cubic 表示の相対同型
\(\mathfrak U_f\simeq[\mathscr G/G_f]\)
（命題~\ref{prop:qa-global-cubic-four-qubit-good-moduli}）
により，矢群のdualとcyclic順序を合わせて
\[
 \boxed{\quad
 [\mathscr G/G_8]\ \longrightarrow\
 [\mathscr G/(G_8/K_{\rm QP})]\simeq\mathfrak U_f
 \quad}
 \tag{HC}\label{eq:qa-cubic-qp-h2-cubic-gerbe}
\]
はbanded $\Gm^3$-gerbeである．
同値は実際のcubic表示の相対groupoidを含むが，
選んだ正規窓・rulingの忘却や任意AS族のheartとの同一視は
追加で検証を要する．
\end{corollary}
\begin{proof}
四頂点のscalar $(\lambda_0,\ldots,\lambda_3)$は
隣接Homに
$(\lambda_1/\lambda_0,\lambda_2/\lambda_1,
\lambda_3/\lambda_2)$で作用する．この三つの値は独立で，
三重tensorの関係部分空間には共通scalarとして作用するから
Grassmann上自明である．
第四矢群のscalarはその逆積となり，表示した
$K_{\rm QP}$に一致する．中心部分群の作用が自明な
商スタックの剛体化について，二つの表示を比べれば
群スキームの完全列が従う．
$K_{\rm vert}$の群作用を割った有効群は
$\operatorname{PGL}_2^3$である．

幾何条件は四重テンソルの不変開条件である．
もし第四矢群に関するflatteningの階数が一以下なら，
四因子テンソルを二因子対二因子の$4\times4$行列として見たとき
その階数は高々二となる．
しかしgeometric条件はその行列と転置のkernelがSegreを避ける
ことを含み，階数二以下を許さない
（補題~\ref{lem:qa-cubic-adjugate-two-invariants}）．
従って$\mathscr G\subset P_8^{(2)}$である．
$f$可逆な$cubic$ H2 chartにおける残余群$G_f$は，
前の四重テンソル表示の四つの矢群の有効群であり，
双対に使う因子についてscalarを逆に読めば
$G_f\simeq G_8/K_{\rm QP}$となる．
命題~\ref{prop:qa-global-cubic-four-qubit-good-moduli}は
この商を実際の相対cubic H2 chartに同定するので，
結論を得る．
\end{proof}

\begin{corollary}[二cubic正規化を貼る操作は一般には右同値でない]
\label{cor:qa-cubic-qp-inversion-not-right-equivalent}
両cubic chartの交わり$\Omega^\times$上で、
$f$可逆な正規化と$g$可逆な正規化から得る二つの
四重テンソル$T,T^{-1}$は、一般点で頂点を固定する
八矢weight-preserving right-equivalenceでは結ばれない。
同値となり得る点の粗good moduli像は
\[
 e_3=de_1\quad\text{または}\quad e_3=-de_1
 \tag{Fix}\label{eq:qa-cubic-qp-inversion-right-equivalence-locus}
\]
という二つの固有閉部分の和に含まれる。
従って二cubic chartの全域比較には、
八矢QPの単一のright-equivalence商だけでなく、
実際のmutation・window変更に対応する貼合せが必要である。
\end{corollary}
\begin{proof}
命題~\ref{prop:qa-cubic-weighted-inversion-gluing}によって、
二つのchartの四重テンソル不変式は
$[e_1:e_2:d:e_3]$と
$\iota[e_1:e_2:d:e_3]=[e_3:d^2e_2:d^3:d^4e_1]$
である。
weight-preserving right-equivalenceなら四因子の
$\GL_2^4$作用による同じorbitであるため、その四量子ビットGIT
同次不変式も一致する。
命題~\ref{prop:qa-cubic-inversion-weight-two-fixed-locus}により、
この等式の固定部分は正に表示した二つの閉部分である。
それらが真であることは、例えばCartanの固有値を
$(1,2,3,4)$とした可逆行列の不変式
\[
 (e_1,e_2,d,e_3)=(30,273,24,820)
\]
からも従う。この行列では
$p_{\rm adj}=9625$、$q_{\rm adj}=6510100$が非零なので
$\Omega^\times$に属し、一方$820\ne\pm24\cdot30$である。
従って一般点で右同値にはならない。
\end{proof}

\begin{lemma}[nilpotent mass上の対角慣性は非平坦になり得る]
\label{lem:qa-cubic-qp-nilpotent-mass-scaling}
一組の重み二の可縮pair $x,y$を八矢potential $H$に加えて
\[
 W_\lambda=H+\lambda yx
\]
とする。$x\mapsto tx,\ y\mapsto sy$で与えられる
二つの矢の対角スカラー変換のうち、
potentialを固定する部分群スキームは
\[
 \mathcal T_\lambda=
 \operatorname{Spec}
 \frac{R[t^{\pm1},s^{\pm1}]}{(\lambda(ts-1))}
 \ \subset\ (\mathbb G_{m,R})^2.
 \tag{Mass}\label{eq:qa-qp-mass-scaling-group}
\]
$\lambda$が可逆なら$\mathcal T_\lambda\simeq\Gm_R$、
$\lambda=0$なら$\mathcal T_\lambda\simeq(\Gm_R)^2$である。
特に$R=k[\epsilon]/(\epsilon^2)$、$\lambda=\epsilon$では
$\mathcal T_\lambda$は$R$上平坦でない。

この補題は通常のright-equivalence商内の明示的な部分慣性の
計算であり、$W_\lambda$が全次数flatなAS正則族をなすこと、
あるいは全慣性群スキームが$\mathcal T_\lambda$に等しいことは
主張しない。
\end{lemma}
\begin{proof}
$H$は$x,y$を含まず、変換によって
$W_\lambda$は$H+\lambda st\,yx$へ移る。
二種類の巡回道は独立だから固定条件は
$\lambda(st-1)=0$に等しい。
これは任意可換$R$-代数への基底変更後も同じ等式である。
$\lambda$が単元なら$st=1$、
$\lambda=0$なら条件は空である。

双対数環での非平坦性を証明する。
表示した環を$A$とし、$z=st-1\in A$と置く。
$\epsilon z=0$だが、$z$は$A/\epsilon A
=k[t^{\pm1},s^{\pm1}]$で非零である。
もし$A$が$R$上flatなら、$R$の完全列
$0\to(\epsilon)\to R\to k\to0$
を$A$でtensorすることで
$\ker(\epsilon:A\to A)=\epsilon A$となる。
これは$z\notin\epsilon A$と矛盾する。
\end{proof}


\subsection{非最小十矢のmassゲージと相対群oid}
\label{subsec:qa-ten-arrow-mass-gerbe}

次は八矢領域から階数三のscheme壁への移行を扱う．
四辺のrank-two矢束を$U,V,W,T$とし，偶数対角の
$x:0\to2,\ y:2\to0$を追加する．
$Q_8$の矢だけで作る道
$q:0\to2$，$p:2\to0$を用いて，
\[
 \mathcal W_\lambda=H+yq+px+\lambda yx
 \tag{MG}\label{eq:qa-mass-gauge-potential}
\]
とする．ここで$H$は四辺を一周するquartic potentialである．
$q\in V\otimes U$，$p\in T\otimes W$であり，これらを
$x,y$に依存しないものとして保つ．
$\lambda$は任意の可換係数環$R$の元で，零因子でもよい．

\begin{theorem}[massを横断する滑らかなゲージ群]
\label{thm:qa-mass-gauge-smooth-group}
\[
 \mathcal U_\lambda
 =\Spec R[t,(1+\lambda t)^{-1}],\qquad
 t\star s=t+s+\lambda ts
 \tag{UG}\label{eq:qa-mass-gauge-group-law}
\]
は$R$上smoothかつflatな相対次元一の可換群スキームである．
その$\mathcal W_\lambda$への重みを保つ右同値作用は
$a=1+\lambda t$として
\[
 x\longmapsto ax+tq,\qquad
 y\longmapsto a^{-1}(y-tp),
 \qquad U,V,W,T\text{の各矢は固定}
 \tag{GA}\label{eq:qa-mass-gauge-action}
\]
で与えられる．この右同値はpotentialを巡回同値の意味で
厳密に固定し，$\Jac(Q_{10},\mathcal W_\lambda)$には
恒等自己同型を誘導する．
$R/(\lambda)$上では$\mathcal U_0\simeq\mathbb G_a$，
$\lambda$が単元の開では
$t\mapsto1+\lambda t$で$\mathcal U_\lambda\simeq\Gm$となる．
$\lambda=\epsilon$ on $R=k[\epsilon]/(\epsilon^2)$でも
$\mathcal U_\epsilon$はsmooth/flatである．
$(\lambda,q)$が係数環上単位idealを生成するとき，この作用は
right-equivalenceの安定化群への群スキーム単射である．
\end{theorem}
\begin{proof}
$1+\lambda t$を可逆化した$\mathbb A^1_R$の主開なので
smooth/flatは直ちに従う．積の式に対し
$1+\lambda(t\star s)=(1+\lambda t)(1+\lambda s)$，
逆元は$-t/(1+\lambda t)$である．
右同値の逆写像も同じ形式で与えられる．

式~\eqref{eq:qa-mass-gauge-action}をpotentialに代入する．
$x'=ax+tq$，$y'=a^{-1}y-ta^{-1}p$と置くと，
$y'q,px',\lambda y'x'$の展開で，
$yx,yq,px,pq$の巡回係数はそれぞれ
\[
 \lambda,\qquad
 a^{-1}(1+\lambda t)=1,\qquad
 a(1-\lambda t/a)=1,\qquad
 t-t/a-\lambda t^2/a=0 .
\]
$pq$と$qp$は巡回同値であり，ゆえに
$\mathcal W_\lambda(x',y')\sim_{\rm cyc}
\mathcal W_\lambda(x,y)$である．
さらに
\[
 x'-x=t(q+\lambda x)=t\,\partial_y\mathcal W_\lambda,
 \qquad
 y'-y=-\frac{t}{a}(p+\lambda y)
      =-\frac{t}{a}\,\partial_x\mathcal W_\lambda .
\]
従ってJacobian商では各生成矢の像が等しい．
$\lambda$が単元なら$a$が任意の可逆元を走り，
$\lambda=0$なら加法群となる．
もし$\lambda$が単元でない点で$q$の一成分が単元なら，
$x$に加えた$q$の係数から$t$が復元される．
$\lambda$が単元なら$x$のスカラー$a$から$t$が復元される．
単位ideal条件によりこれらがZariski局所に適用でき，
作用は単射である．
\end{proof}

\begin{theorem}[一積可逆なscheme壁の十矢QP群oid]
\label{thm:qa-ten-arrow-relative-finite-gerbe}
$Z^{\rm nd}$と$G$を有限foundationの
定義~\ref{def:qa-foundation-open}・命題~\ref{prop:qa-foundation-stack}
に従うものとする．$b:W\otimes V\to F$が可逆な開の
unrigidified stackを
\[
 \mathfrak F^{\rm lin}_{b\text{-inv}}
   =[(Z^{\rm nd}\cap D(\det b))/G]
\]
と置く．十矢potential空間のうち
式~\eqref{eq:qa-mass-gauge-potential}の係数が次の
二条件を満たす開groupoidを$\mathfrak P^{\rm good}_{10,e}$
とする：
\begin{enumerate}
 \item 全幾何点で，$\lambda\ne0$または$q$が非分解可能
       （すなわち$D(\lambda)\cup D(\det q)$）．
 \item $\mathcal D=(V\otimes U)\oplus\mathcal O x$，
       $r=q+\lambda x$とし，
       $\partial_T\mathcal W_\lambda:T^\vee\to W\otimes\mathcal D$
       を取ると，次の写像の像が階数四の直和因子である：
       \[
       (W\otimes\mathcal O r)\oplus T^\vee
          \longrightarrow W\otimes\mathcal D,\qquad
       (w\otimes r,\tau)\longmapsto
             w\otimes r+\partial_T\mathcal W_\lambda(\tau).
       \tag{FR}\label{eq:qa-mass-gauge-relation-bundle}
       \]
\end{enumerate}
この開はweighted right-equivalence groupoidの制限として
任意の基底変更に安定である．
その各対象から
\[
 E=\mathcal D/(\mathcal O r),\quad
 F=W\otimes V,\quad
 N=(W\otimes\mathcal D)/
       (W\otimes r+\operatorname{im}\partial_T\mathcal W_\lambda)
 \tag{FD}\label{eq:qa-mass-gauge-finite-reconstruction}
\]
を構成すると，有限foundationの全積を保存する射
\[
 \Phi:\mathfrak P^{\rm good}_{10,e}
          \longrightarrow\mathfrak F^{\rm lin}_{b\text{-inv}}
\]
が得られる．$\Phi$はfppf局所にessentially surjectiveで，
全Isom sheafの射もfppf局所に全射である．
各対象における自己同型の核は
定理~\ref{thm:qa-mass-gauge-smooth-group}の
$\mathcal U_\lambda$に正確に等しい．
すなわち$\Phi$はfppf局所に$\mathcal U_\lambda$を
相対慣性とするgerbeであり，幾何的rank三の壁上では
$\mathbb G_a$，可逆cubic開では$\Gm$となる．
$\lambda$がnilpotentでも群スキームはsmooth/flatである．
この定理は有限代数のgroupoidについての主張であり，
任意のpotentialのAS正則性を含まない．
\end{theorem}
\begin{proof}
第一条件は幾何点での条件であり，
$\lambda=0$では$q$の非分解性は二因子の基底変更と
可逆scalar倍で保存される．
$x$への二段道の付加は$q$を$\lambda$倍の道だけ変えるので，
この条件は右同値に不変である．第二条件も，関係束とその
直和因子性の不変性から保たれる．

$r$は各幾何点で非零のprimitive元であり，
$\mathcal O r\subset\mathcal D$はrank-one直和因子である．
第二条件によって$E,N$はそれぞれrank四・六の束となる．
二段の積$a:V\otimes U\to E$は商写像，$b$は恒等，
三段の積$c:W\otimes E\to N$は商写像，
$d:(W\otimes V)\otimes U\to N$はその合成である．
従って結合律と$(c,d)$の全射性を満たす．
$\lambda\ne0$なら$a$は可逆であり，
$\lambda=0$では$\ker a=kq$がSegreを避ける．
ゆえに全ての幾何点で$Z^{\rm nd}$の条件を満たす．
$x,y,T$のright-equivalenceはrelationの枠やliftを変更する
だけなので，この構成は相対groupoid射となる．

逆方向を局所的に構成する．
$b$が可逆なfinite foundationでは結合律により
$d$は$c$と$a$から決まり，$(c,d)$の全射性は$c$の全射性と同値である．
$a$は各幾何点でrank三または四である．
ある三次minorを可逆化し，源$V\otimes U$の第四基底を$q$，
標的$E$の第四補元を$x$と選べば
$a\sim\operatorname{diag}(I_3,-\lambda)$となり，
$E\simeq((V\otimes U)\oplus\mathcal O x)/(q+\lambda x)$を得る．
rank三の点では$q$は既存の非分解可能kernel線の生成元なので，
$\det q$は単元に取れる．
$c$の核$K\subset W\otimes E$はrank二の直和因子である．
その局所基底を$T^\vee$で標識し，
$W\otimes\mathcal D\to W\otimes E$へのliftを選べば
$\partial_T\mathcal W_\lambda$の二本を得る．
各liftを三段道成分と$x$成分に分けると
$H+px$が復元し，$y(q+\lambda x)$を追加して
十矢potentialを得る．構成~\eqref{eq:qa-mass-gauge-finite-reconstruction}
は元の全積を回収する．
これらの枠・splitting・liftはfppf局所に取れ，
任意の有限foundation同型も同じ選択を合わせることで
十矢right-equivalenceへ局所に持ち上がる．

最後に自己同型核を計算する．
有限foundationを恒等に保つ右同値は$U,V,W$上で恒等であり，
$x\mapsto\alpha x+z$について
$\alpha x+z-x$が$E$で零である．
$r=q+\lambda x$がprimitiveだから，これは唯一の$t$について
\[
 \alpha=1+\lambda t,\qquad z=tq
\]
と同値である．$\alpha$は単元なので$t\in\mathcal U_\lambda$．
また$\partial_T\mathcal W_\lambda$の$W\otimes E$への像は
rank二直和因子であり，$N$を恒等に固定する
potential自己同型は$T$の基底を変更できない．
残る$y\mapsto\beta y+h$をpotentialの係数と比較する．
$r$がprimitiveであることから
\[
 \beta=(1+\lambda t)^{-1},
 \qquad h=-t(1+\lambda t)^{-1}p
\]
が一意に強制される．これは
式~\eqref{eq:qa-mass-gauge-action}そのものである．
従って核の群スキームは$\mathcal U_\lambda$に尽き，
任意基底変更でも同じ計算が成立する．
\end{proof}

\begin{corollary}[実際のmutation AS族への適用範囲]
\label{cor:qa-mass-gauge-actual-as-window}
定理~\ref{thm:qa-relative-mutation-nonminimal-as-resolutions}
の局所十矢表示
$\mathcal W_\delta=\delta yx+yq+px+H$は
上記$\Phi$の対象である．特に
$\delta=\epsilon$ on $k[\epsilon]/(\epsilon^2)$でも
finite foundationの全成分と全積は平坦に回収され，
raw QP自己同型からfinite foundation自己同型への核には
滑らかな$\mathcal U_\epsilon$が現れる．
補題~\ref{lem:qa-cubic-qp-nilpotent-mass-scaling}の
非平坦な対角scalar部分慣性$\mathcal T_\epsilon$は，
別のraw potential $H+\epsilon yx$ に対する模型である．
実際のmutation AS potential $H+yq+px+\epsilon yx$ に
そのまま同じ部分群があるとは主張しない．
これに対し本定理の$\mathcal U_\epsilon$は
実際のmutation AS族に作用する滑らかなゲージ核である．
全相対heartと選択したmutation窓の一致は
定理~\ref{thm:qa-relative-rot-mut-completion-heart}による．
ここから選択した窓を忘れた全AS/QP stackと
内在的$H_2$の同値を推論しない．
\end{corollary}
\begin{proof}
補題~\ref{lem:qa-relative-mutation-seven-three-presentation}は
$q$のprimitive性と$\delta$の任意の冪零値を許し，
命題~\ref{prop:qa-relative-mutation-ordinary-jacobian}は
同じpotentialが実際のflatな相対有限代数・全次数完成を
与えることを示す．前定理を適用すれば結論が従う．
\end{proof}

\begin{remark}[十二矢QPの余剰慣性と残る比較]
\label{rem:qa-cubic-qp-remaining-inertia}
命題~\ref{prop:qa-stabilized-inertia}で既証明の
\(
 \mathcal P_{12}^{\lambda\mu\ne0}
 \simeq\mathcal P_8\times B(\Gm^2)
\)
をgeometric開に制限すれば、十二矢QPから上の有限
foundationへの忘却には、通常の四頂点scalarの$\Gm^3$に加え
二つの可縮pairの$\Gm^2$がある．
従ってraw十二矢QPの全right-equivalence群oidを
$\mathfrak U_f$へ忠実に同定してはならない．
本項は八矢の相対同型層を可逆cubic開で決着させるものであり，
rank-three壁・noncubic・nilpotent massを持つ十／十二矢での
相対QP群oidや同じheartを保つ比較までは証明しない．
特に\(\operatorname{coker}\mu_G=\OO/(\delta)\)が
非平坦となる一次変形
（命題~\ref{prop:qa-relative-as10-minimal-arrow-wall}）を
幾何fiberごとの最小矢数だけから復元することはできない．
\end{remark}
